\begin{center}
\begin{tikzpicture}[scale=1.1,>=latex]
%-------------------------------------- Defino estilo de flecha -----------------------------------
\tikzset{
    mid arrow/.style={
        decoration={
            markings,
            mark=at position 0.5 with {\arrow{>}} 
        },
        postaction={decorate} 
    },
    init arrow/.style={
        decoration={
            markings,
            mark=at position 0.2 with {\arrow{>}} 
        },
        postaction={decorate} 
    }
}

\begin{axis}[
    hide axis, clip=false,
    view={90}{45}
]

%------------------------------------- Cilindro -------------------------------------
\addplot3 [
    surf, shader=faceted interp,
    point meta=x,
    colormap={graywhite}{color=(gray!70!white) color=(white)}, % Gradiente suave
    samples=40,
    samples y=5,
    z buffer=sort,
    domain=0:360,
    y domain=-0.5:0.5
] (
    {cos(x)},
    {sin(x)},
    {0.5*y}
);

%------------------------------------- Circunferencia central e_3 ----------------------------------------
\addplot3 [
    samples=50,
    domain=85:-85,
    samples y=0,
    thick,
    violet,
    mid arrow
] (
    {cos(x)},
    {sin(x)},
    {0}
) node[pos=0.5, below] {$e_3$}; % <- El nombre se añade aquí
\addplot3 [
    samples=50,
    domain=-115:-245,
    samples y=0,
    thick,
    violet,
] (
    {cos(x)},
    {sin(x)},
    {0}
);

%------------------------------- Corte transversal vertical superior e_1' ------------------------------------
\addplot3 [
    samples=50,
    domain=0:0.5,
    samples y=0,
    thick,
    yellow,
    mid arrow,
] (
    {(1+0.5*x*cos(180/2)))*cos(180)}, % x: sustituye phi = 180
    {(1+0.5*x*cos(180/2)))*sin(180)}, % y: sustituye phi = 180
    {0.5*x*sin(180/2)}                % z: sustituye phi = 180
)node[pos=0.5, left] {$e_1'$}; % <- El nombre se añade aquí

%------------------------------- Corte transversal vertical inferior e_2' ------------------------------------
\addplot3 [
    samples=50,
    domain=-0.5:0,
    samples y=0,
    thick,
    yellow,
    mid arrow,
] (
    {(1+0.5*x*cos(180/2)))*cos(180)}, % x: sustituye phi = 180
    {(1+0.5*x*cos(180/2)))*sin(180)}, % y: sustituye phi = 180
    {0.5*x*sin(180/2)}                % z: sustituye phi = 180
)node[pos=0.5, left] {$e_2'$}; % <- El nombre se añade aquí

%------------------------------------- Circunferencia central e_5' ----------------------------------------
\addplot3 [
    samples=50,
    domain=-85:85,
    samples y=0,
    thick,
    cyan,
    mid arrow
] (
    {cos(x)},
    {sin(x)},
    {-0.25}
) node[pos=0.5, below] {$e_5'$}; % <- El nombre se añade aquí
\addplot3 [
    samples=50,
    domain=-147:-213,
    samples y=0,
    thick,
    cyan,
] (
    {cos(x)},
    {sin(x)},
    {-0.25}
);

%------------------------------------- Circunferencia superior e_4' ----------------------------------------
\addplot3 [
    samples=50,
    domain=-180:180,
    samples y=0,
    thick,
    red!70!yellow,
    mid arrow,
] (
    {cos(x)},
    {sin(x)},
    {0.25}
)node[pos=0.5, below] {$e_4'$}; % <- El nombre se añade aquí

%---------------------------------------------- Vértices ----------------------------------------------
% Vértices en phi = 180°
\node[right, yshift=7pt] at (axis cs:-1,0,-0.25) {$v_3'$};
\node[right, yshift=7pt] at (axis cs:-1,0,0) {$v_2$};
\node[right, yshift=7pt] at (axis cs:-1,0,0.25) {$v_1'$};
\addplot3[only marks, mark=*, mark size=1.5pt, black] coordinates {(-1,0,-0.25) (-1,0,0) (-1,0,0.25)};
\end{axis}
\end{tikzpicture}
\end{center}
